\subsection{Isolated points and sensitive points of the spectrum}

Let E be a complete normable space and let u be an endomorphism of E. Let \(\lambda\) be an isolated point of the spectrum of u. Recall that we denote by \(e_{\lambda}(u)\) the spectral projector associated with u and the closed and open set \(\{\lambda\}\) of the spectrum of u (no. 3 of I, p. 131). The endomorphism \(u - \lambda1_{E}\) induces an automorphism of \(\operatorname{Ker}(e_{\lambda}(u))\) and a quasi-nilpotent endomorphism of \(\operatorname{Im}(e_{\lambda}(u))\) (loc. cit.).

DEFINITION 1. --- One says that \(\lambda\) has finite spectral multiplicity for \(u\) if the spectral projector \(e_{\lambda}(u)\) is of finite rank; in this case, the integer \(m_{\lambda}(u) = \dim (\mathrm{Im}(e_{\lambda}(u)))\) is called the spectral multiplicity of \(\lambda\) for \(u\) .

DEFINITION 2. --- One calls sensitive points of the spectrum of u the isolated points of Sp(u) of finite spectral multiplicity. The set of sensitive points of Sp(u) is called the sensitive spectrum of u and is denoted Sp \(_{s}\) (u). \(^{(2)}\) One calls essential spectrum of u and denotes by Sp \(_{e}\) (u) the complementary set \(\mathrm{Sp}(u) - \mathrm{Sp}_{\mathrm{s}}(u)\).

As the set \(\mathrm{Sp}_{\mathrm{s}}(u)\) consists of isolated points of \(\mathrm{Sp}(u)\) , it is open in \(\mathrm{Sp}(u)\) , discrete, and countable (III, p. 78, lemma 3); it is also bounded. The set \(\mathrm{Sp}_{\mathrm{e}}(u)\) is closed in \(\mathrm{Sp}(u)\) , hence compact.

By prop. 14 of III, p. 54, the sensitive points of the spectrum of \(u\) are the complex numbers \(\lambda\) such that \(u - \lambda 1_{\mathrm{E}}\) is a Riesz endomorphism of E that is not an automorphism of E.

For every complex number \(\lambda\) , denote by \(\mathrm{N}_{\lambda}(u)\) and \(\mathrm{I}_{\lambda}(u)\) the nilspace and the conilspace of the endomorphism \(u - \lambda 1_{\mathrm{E}}\) of E. Since E is a Fréchet space, \(u - \lambda 1_{\mathrm{E}}\) is a Riesz endomorphism if and only if \(\mathrm{N}_{\lambda}(u)\) is finite-dimensional and \(\mathrm{I}_{\lambda}(u)\) has finite codimension (prop. 13 of III, p. 53); when this is the case, \(u - \lambda 1_{\mathrm{E}}\) is invertible if and only if \(\mathrm{N}_{\lambda}(u)\) is reduced to 0 (III, p. 47, prop. 6). The sensitive points of the spectrum of \(u\) are therefore the complex numbers \(\lambda\) such that \(\mathrm{N}_{\lambda}(u)\) is of nonzero finite dimension and \(\mathrm{I}_{\lambda}(u)\) is of finite codimension in E (III, p. 54, prop. 14). The spectral multiplicity \(m_{\lambda}(u)\) of \(\lambda\) is then the dimension of \(\mathrm{N}_{\lambda}(u)\) , and the endomorphism \(u - \lambda 1_{\mathrm{E}}\) induces by restriction an automorphism of \(\mathrm{I}_{\lambda}(u)\) and a nilpotent endomorphism of \(\mathrm{N}_{\lambda}(u)\) . Since \(\mathrm{N}_{\lambda}(u)\) is not reduced to 0, it follows that \(\lambda\) is an eigenvalue of \(u\) . The smallest integer \(n \geqslant 1\) such that \(\operatorname{Ker}((u - \lambda 1_{\mathrm{E}})^n)\) is equal to \(\mathrm{N}_{\lambda}(u)\) is the order of the pole of the resolvent of \(u\) at the point \(\lambda\) (No. 3 of I, p. 131), which is bounded above by \(m_{\lambda}(u)\) .

In particular, we have therefore proved:

PROPOSITION 1. --- The sensitive points of the spectrum of u are eigenvalues of finite spectral multiplicity.

The essential spectrum of \(u\) consists of the complex numbers \(\lambda\) such that \(u - \lambda 1_{\mathrm{E}}\) is not a Riesz endomorphism of E. In particular, if \(u\) is compact, then every \(\lambda \in \mathrm{Sp}(u) - \{0\}\) is a sensitive point of the spectrum (cor. 2 of III, p. 77).

PROPOSITION 2. --- Let E be a complete normable space and u an endomorphism of E. Let H be a finite subset of \(\mathrm{Sp}_{\mathrm{s}}(u)\) . The set H is open and closed in \(\mathrm{Sp}(u)\) . The spectral projector \(e_{\mathrm{H}}(u)\) has finite rank, has image \(\sum_{\lambda \in \mathrm{H}}\mathrm{N}_{\lambda}(u)\) and kernel \(\bigcap_{\lambda \in \mathrm{H}}\mathrm{I}_{\lambda}(u)\) .

This follows from no. 3 of I, p. 131 since, for every \(\lambda \in \mathrm{Sp}_{\mathrm{s}}(u)\) , the subspaces \(\mathrm{N}_{\lambda}(u)\) and \(\mathrm{I}_{\lambda}(u)\) are respectively the image and the kernel of the spectral projector \(e_{\lambda}(u)\) .

COROLLARY 1. --- Let V be a neighborhood of \(\mathrm{Sp_e}(u)\) .

\begin{enumerate}
\def\labelenumi{\alph{enumi})}
\item
  There exists a decomposition of E as a topological direct sum F ⊕ G such that F is finite-dimensional, F and G are stable under every endomorphism commuting with u, and the spectrum of the endomorphism of G induced by u is contained in V.
\item
  Suppose V is nonempty. There exists an element v of the bicommutant of u in \(\mathcal{L}(\mathrm{E})\) whose spectrum is contained in V and such that \(v - u\) has finite rank.
\end{enumerate}

Set \(\mathrm{H} = \mathrm{Sp}(u) \cap (\mathbf{C} - \mathring{\mathrm{V}})\) . The set H is contained in \(\mathrm{Sp}_{\mathrm{s}}(u)\) , hence is discrete; since it is closed in the compact space \(\mathrm{Sp}(u)\) , it is finite. One may take for F and G the image and the kernel of the spectral projector \(e_{\mathrm{H}}(u)\) (prop. 2).

Suppose that V is not empty. Let \(\mu\in\mathrm{V}\) and denote by \(v\) the endomorphism of E that coincides with the homothety of ratio \(\mu\) in F and with \(u\) in G. Its spectrum is contained in \(\{\mu\}\cup(\mathrm{Sp}(u)-\mathrm{H})\) , therefore in V, and \(v-u\) is of finite rank since F is finite-dimensional.

For every compact subset S of C, we denote by \(\mathcal{O}(S)\) the algebra of germs of holomorphic functions in a neighborhood of S and with values in C (I, p. 49, §4, no. 1).

COROLLARY 2. --- Let \(f \in \mathcal{O}(\mathrm{Sp}(u))\) and let \(\mu\) be a complex number. Let H be the set of \(\lambda \in \mathrm{Sp}(u)\) such that \(f(\lambda) = \mu\) . Then \(\mu\) is a sensitive point of the spectrum of \(f(u)\) if and only if H is finite, nonempty, and contained in the sensitive spectrum of \(u\) . Under these conditions, the spectral projector \(e_{\mu}(f(u))\) is equal to the projector \(e_{\mathrm{H}}(u)\) associated with \(u\) and H. We have

\[
\mathrm{N} _ {\mu} (f (u)) = \bigoplus_ {\lambda \in \mathrm{H}} \mathrm{N} _ {\lambda} (u), \quad \mathrm{I} _ {\mu} (f (u)) = \bigcap_ {\lambda \in \mathrm{H}} \mathrm{I} _ {\lambda} (u)
\]

and the spectral multiplicity of \(\mu\) for \(f(u)\) is the sum of the spectral multiplicities of the elements of H, that is,

\[
m _ {\mu} (f (u)) = \sum_ {\lambda \in \mathrm{H}} m _ {\lambda} (u).
\]

The complex number \(\mu\) belongs to \(f(\mathrm{Sp}(u))\) , and therefore to the spectrum of \(f(u)\) (prop. 8 of I, p. 75), if and only if H is nonempty. We then have \(e_{\mu}(f(u)) = e_{\mathrm{H}}(u)\) (no. 1 of I, p. 127). The other assertions follow from this by Prop. 2.

Examples. --- 1) Let E be a finite-dimensional vector space and u an endomorphism of E. The spectrum of u is finite and coincides with \(\mathrm{Sp}_{\mathrm{s}}(u)\) . Its elements are the roots of the characteristic polynomial \(\chi_{u}\) of u; these are the eigenvalues of u. The spectral multiplicity \(m_{\lambda}(u)\) of an element \(\lambda \in \mathrm{Sp}(u)\) is the multiplicity of \(\lambda\) as a root of the polynomial \(\chi_{u}\) (A, VII, p. 36, cor.). By cor. 2 above, we have

\[
\chi_ {f (u)} (\mathrm{T}) = \prod_ {\lambda \in \operatorname{Sp} (u)} (\mathrm{T} - f (\lambda)) ^ {m _ {\lambda}}
\]

for all \(f \in \mathcal{O}(\mathrm{Sp}(u))\) . This generalizes prop. 10 of A, VII, p. 36.

\begin{enumerate}
\def\labelenumi{\arabic{enumi})}
\setcounter{enumi}{1}
\tightlist
\item
  Let X be a compact subset of C and \(\mathcal{C}(X)\) the Banach space of continuous functions on X with complex values. Denote by u the endomorphism of \(\mathcal{C}(X)\) that sends \(f \in \mathcal{C}(X)\) to the function \(z \mapsto zf(z)\) in \(\mathcal{C}(X)\) . The spectrum of u is equal to X, its sensitive points are the isolated points of X and their spectral multiplicities are equal to 1.
\end{enumerate}

